\section{AI You Can Actually Trust}

Every platform in the category is now shipping AI. The question that matters is not
\emph{whether} a vendor has an LLM feature --- it is \emph{where the model sits
relative to your numbers}. Bolt a generative model onto a black-box ledger and you
have added a confident, occasionally-wrong author with write access to financial
truth.

Ratio's architecture answers the question structurally. Large language models live
\emph{above} a hard trust boundary; the proven kernel lives below it. A model can:

\begin{itemize}
  \item translate an advisor's plain-language intent into a precise, typed
  specification;
  \item draft client-ready explanations of a position, a fee, or a performance
  result; and
  \item propose posting rules, mappings, and reconciliations for human review.
\end{itemize}

% Scoped to conservation: the kernel proves conservation, not authorization.
% Authorization/access control is a separate control plane, not a kernel
% theorem, so we do not claim it here (crank-001 finding C1/D2).
\noindent What a model \emph{cannot} do is write an entry that violates
conservation. Every proposal is compiled and checked against the kernel's
conservation invariant before it can affect a ledger; anything that fails to
conserve value is rejected at the boundary. The worst case for a hallucination
is a discarded suggestion --- never an unbalanced book.

\begin{callout}
\sffamily\color{ratioInk}
\textbf{Net:} you get the productivity of an AI copilot for the parts that are
genuinely language problems --- intent, explanation, drafting --- with a
mathematical guarantee that the model can never be the reason a number is wrong.
\end{callout}
